\documentclass{article}
\usepackage{amssymb, latexsym, amsmath, graphics, fullpage, epsfig, amsthm, relsize, pgf, tikz, amsfonts, makeidx, latexsym, ifthen, hyperref, calc}
%\usepackage{eucal} this is a different font than \mathcal{•}
\usetikzlibrary{arrows}
\newtheorem{theorem}{Theorem}[section]
\newtheorem{remark}{Remark}[section]
\newtheorem{corollary}{Corollary}[section]
\newtheorem{lemma}{Lemma}[section]
\newtheorem{assumptions}{Assumptions}[section]
\newtheorem{proposition}{Proposition}
\newtheorem{definition}{Definition}
\newtheorem{notation}{Notation}
\usepackage{mathrsfs}
\usepackage[shortlabels]{enumitem}
\usepackage{tikz}
\usepackage{amsmath}
\usetikzlibrary{arrows}
\usetikzlibrary{shadows.blur}
\usepackage{pgfplots}
\usepackage{mwe} % For dummy images
\usepackage{subcaption}
\usepackage{multirow}
\usepackage{dcolumn}
\newcolumntype{2}{D{.}{}{2.0}}
\usepackage{multicol}
\numberwithin{equation}{section}

\usepackage[notref,notcite]{showkeys}

\newcommand{\lip}{\textup{Lip}_b}
\newcommand{\liph}{\text{Lip}_{\hat\rho}}
\newcommand{\R}{\mathbb{R}}
\newcommand{\E}{\mathbb{E}}
\newcommand{\Norm}[1]{\left\|  #1   \right\|}
\newcommand{\ud}{\ensuremath{\mathrm{d}}}

\usepackage{lipsum}%% a garbage package you don't need except to create examples.
\usepackage{fancyhdr}
\pagestyle{fancy}
%\lhead{\Huge Version A}
%\rhead{\thepage}
\cfoot{}
\renewcommand{\headrulewidth}{0pt}

\usepackage[headheight = .5in, headsep = \baselineskip, top = 1in, left = 1in, right = 1in, textwidth = 7.52 in]{geometry}
\lhead{ \parbox[][\headheight][t]{5cm}{\textbf{}}}
\rhead{\parbox[][\headheight][t]{2cm}{\raggedleft Page\,\thepage{} of \pageref{LastPage}}}
\usepackage{lastpage}

\usepackage{tikz}
\usepackage{pgfplots}
\pgfplotsset{compat=newest}
\usetikzlibrary{decorations.markings}

\allowdisplaybreaks


\begin{document}

\section{(2.11) verification}
The admissible weight property
\[
\left(G(t,\cdot)*\rho(\cdot)\right)(x) =(2\pi t)^{-d/2} \int_{\mathbb R^d}\exp\left(-\frac{|x-y|^2}{2t}\right)\rho(y)dy \le K_T\rho(x) \quad \forall x \in \mathbb R^d,\; K_T \in \mathbb N,\; t\in[0,T]
\]
tells us that
\begin{equation} \label{aw}
\lim_{|x| \to \infty}\frac{\left(G(t,\cdot)*\rho(\cdot)\right)(x)}{\rho(x)} \le K_T \quad t \in [0,T].
\end{equation}

\bigskip
\textcolor{magenta}{The expression for $K_t$ should be
\begin{align*}
	K_t:= \sup_{x\in\R^d} \frac{\left(G(t,\cdot)*\rho(\cdot)\right)(x)}{\rho(x)}
\end{align*}
The limit in \eqref{aw} gives a sufficient (also necessary) condition for $K_t$ to be finite, but in general, this limit is not equal to the value of $K_t$.
}

\textcolor{blue}{I meant for $K_T$ to mean $C_\rho(T)$. I made some corrections above}


We want to show that
\begin{equation} \label{con}
\exp\left(-a|x|^2\right) \le K_a'\rho(x)
\end{equation}

\bigskip
\textcolor{magenta}{Why does \eqref{aw} imply \eqref{con}?}

\textcolor{blue}{\eqref{con2} and \eqref{aw} imply \eqref{con1} which implies \eqref{con}}

and it suffices to show that
\begin{equation} \label{con1}
\lim_{|x| \to \infty}\frac{\exp\left(-a|x|^2\right)}{\rho(x)} \le K_a'.
\end{equation}
since then
\[
\exp\left(-a|x|^2\right) \le (K_a'+\epsilon)\rho(x) \quad |x|>R_\epsilon
\]
and by the positivity and continuity of both $	\exp\left(-a|x|^2\right)$ and $\rho(x)$, we know that there exists some positive constant $K(a,\epsilon)$ such that
\[
\exp\left(-a|x|^2\right) \le K'(a,\epsilon)\rho(x) \quad x\le R_\epsilon
\]
and by then letting $K(a,\epsilon) = \max\{K'(a,\epsilon),(K_a'+\epsilon)\}$ then
\[
	\exp\left(-a|x|^2\right) \le K(a,\epsilon)\rho(x) \quad x \in \R^d
\]
which is what we want to show.

In order to show \eqref{con1}, it suffices to show that
\begin{equation} \label{con2}
	\lim_{|x| \to \infty} \frac{\left(G(t,\cdot)*\rho(\cdot)\right)(x)}{\exp\left(-a|x|^2\right)} = \infty
\end{equation}
because this would mean that $\exp\left(-a|x|^2\right)$ is much smaller than $\left(G(t,\cdot)*\rho(\cdot)\right)(x)$ for all large $|x|$ and since \eqref{aw} is true then this would mean that \eqref{con1} would be true as well by squeeze theorem,
\[
	0 \le \lim_{|x| \to \infty}\frac{\exp\left(-a|x|^2\right)}{\rho(x)} \le  \lim_{|x| \to \infty} \frac{\left(G(t,\cdot)*\rho(\cdot)\right)(x)}{\rho(x)} \le K_t
\]

\eqref{con2} is proved as follows: first recall that $|\rho(x)| < M$ for some $M \in \mathbb N$. Then,
\begin{align*}
0 \le	\dfrac{\int_{\mathbb R^d}\exp\left(-\frac{|x-y|^2}{2t}\right)\rho(y)\ud y}{\exp\left(-a|x|^2\right)} &\le \dfrac{M\int_{\mathbb R^d}\exp\left(-\frac{|x-y|^2}{2t}\right) \ud y}{\exp\left(-a|x|^2\right)}
\\&  = \dfrac{M*\sqrt{2 \pi t}}{\exp\left(-a|x|^2\right)}
\\& \to \infty \quad \text{as } |x|\to \infty
\end{align*}

Thus \eqref{con2} is proved which means that \eqref{con} is also true.

\bigskip

\textcolor{magenta}{Can you give an explicit form of $C_\rho(T)$ so that we can check whether $\sup_{T>0}C_\rho(T)<\infty$, or your $K_t$? If not, can you derive some sufficient conditions that guarantee this
property?}

\textcolor{blue}{In the above wrote $K_t$ to mean the same thing as $C_\rho(T)$.
	However in general there is no explicit form for this constant and also in general, $\sup_{T \ge 0}K_T \not < \infty$.
For example, when we consider the weight $\rho(x)=\exp(-|x|)$ in $\R$ then we can show that $(G(t,\cdot)*\exp(-|\cdot|))(x) \le 2\exp(t/2) \exp(-|x|)$ and clearly $\sup_{t \ge 0} 2\exp(t/2) = \infty$.}
\bigskip

\textcolor{magenta}{This is a good example. We should not expect $\sup_{T>0}C_\rho(T)<\infty$. Actually, you may want to check eq. (4.4) of my AOP paper with Robert 2015 on p. 3040. }

\section{Mistake involving the $G(t,x) \le K G(1,x)$ inequality}
During a meeting we had, we proved that $G(t,x) \le K G(1,x)$ where the constant $K$ does not depend on $x$ or $t$ but we made a mistake in doing so. What we did is the following (for now we work in $\R$):

\begin{align*}
\frac{G(t+1,x)}{G(1,x)} &= \sqrt{\frac{1}{t+1}}\exp\left( -\frac{|x|^2}{2(t+1)} + \frac{|x|^2}{2}\right)
\\& = \sqrt{\frac{1}{t+1}}\exp\left( -\frac{|x|^2}{2} \left(-\frac{1}{t+1}+ 1\right)\right) \quad \text{mistake was made here when factoring}
\\& = \sqrt{\frac{1}{t+1}}\exp\left( -\frac{|x|^2}{2}\frac{t}{t+1}\right)
\\& \le K
\end{align*}
However if we fix the mistake, what we should get is that
\begin{align*}
\frac{G(t+1,x)}{G(1,x)} &= \sqrt{\frac{1}{t+1}}\exp\left( -\frac{|x|^2}{2(t+1)} + \frac{|x|^2}{2}\right)
\\& = \sqrt{\frac{1}{t+1}}\exp\left( -\frac{|x|^2}{2} \left(\frac{1}{t+1}- 1\right)\right)
\\& = \sqrt{\frac{1}{t+1}}\exp\left( \frac{|x|^2}{2} \frac{t}{t+1} \right)
\\& \not \le K
\end{align*}


\textcolor{magenta}{Let's remove the property $\sup_{t\ge 0}\E\left(\left|u(t,\cdot\right|_\rho^2\right)<\infty$ unless you can show that $C_\rho(T)$ is bounded in $T$.}

\section{Mistake that $G(t+t_0,x)$ is bounded above by $G(t,x)$ for $t > t_0$}

\begin{align*}
	\frac{G(t+t_0,x)}{G(t_0,x)} &= \sqrt{\frac{t}{t+t_0}}\exp\left(-\frac{|x|^2}{2(t+t_0)}\right)\exp\left(\frac{|x|^2}{2t_0}\right)
	\\& = \sqrt{\frac{t_0}{t+t_0}} \exp\left(\frac{|x|^2}{2(t+t_0)}\right)
\end{align*}
and this is not bounded above by a constant independent of $x$ and $t$.

\bigskip
\textcolor{magenta}{This makes more sense!}

\section{Admissible Weight}
\subsection{$\rho(x) = \exp(-|x|)$}

Let $\rho(x) = \exp(-|x|)$.
\begin{align*}
\sqrt{2\pi t} \; (G(t,\cdot) * \rho)(x) &= \int_{\R} \exp\left(-\frac{|y|^2}{2t} \right) \exp\left( -|x-y| \right) \ud y
\\& \le \int_{\R} \exp\left(-\frac{|y|^2}{2t} \right) \exp\left( -|x| +|y| \right) \ud y \quad \quad \text{since } -|x-y| \le -|x| + |y|
\\& = \exp(-|x|) \int_{\R} \exp\left(-\frac{|y|^2}{2t} \right) \exp\left( |y| \right) \ud y
\\& = \exp(-|x|) \int_{\R} \exp\left(-\frac{|y|^2 - 2t|y|}{2t} \right) \ud y
\\& = \exp(-|x|) \exp(t/2) \int_{\R} \exp\left(-\frac{(|y|-t)^2 }{2t} \right) \ud y
\\& =\exp(-|x|) \exp(t/2)  \sqrt{2\pi t} \left[ \text{erf}\left(\sqrt{\frac{t}{2}}\right) +1 \right]
\end{align*}
and therefore
\[
(G(t,\cdot) * \rho)(x) \le \rho(x)\exp(t/2) \left[ \text{erf}\left(\sqrt{\frac{t}{2}}\right) +1 \right].
\]
Now fix $T>0$. Since $\exp(t/2)\left[ \text{erf}\left(\sqrt{t/2}\right) +1 \right]$ is increasing in $t$, then for every $t\in [0,T]$ we have that
\[
(G(t,\cdot) * \rho)(x) \le \exp(T/2)\left[ \text{erf}\left(\sqrt{\frac{T}{2}}\right) +1 \right]  \rho(x)
\]
which shows that $\rho(x) = \exp(-|x|)$ is admissible with $C_\rho(T) = \exp(T/2)\left[ \text{erf}\left(\sqrt{\frac{T}{2}}\right) +1 \right]$, which is the sharpest such constant for the particular $\rho$ given in this example.

\subsection{$\rho(x) = \exp(-|x|^2)$}

In this case
\[
\rho(x) = \sqrt{\pi} G( 1/2 , x )
\]
and so we see that
\begin{equation}\label{E:exp(x^2)}
(G(t,\cdot) * \rho)(x)  = \sqrt{\pi} (G(t,\cdot) * G(1/2, \cdot)(x) = \sqrt{\pi} G(t+1/2,x) = \frac{1}{\sqrt{2t+1}} \exp\left(\frac{-|x|^2}{2t+1}\right).
\end{equation}
I claim that for any $T>0$, that there does not exist a constant $K(T)$ such that
\[
(G(t,\cdot) * \rho)(x) \le K(T) \exp(-|x|^2) \quad \forall \; t\in[0,T] \; \text{and} \; x\in \R.
\]
BWOC, suppose there did. Then for $T=4$, there exists some constant $K(4) $ such that for any $t\in[0,4]$,
\[
\frac{1}{\sqrt{2t+1}} \exp\left(\frac{-|x|^2}{2t+1}\right) \le K(4) \exp(-|x|^2).
\]
Note that for any  $0\le t \le 4$ and $x\in \R$,
\[
\frac{1}{3} \exp\left(\frac{-|x|^2}{2t+1}\right)  \le \frac{1}{\sqrt{2t+1}} \exp\left(\frac{-|x|^2}{2t+1}\right)
\]
and so we must have
\[
\frac{1}{3} \exp\left(\frac{-|x|^2}{2t+1}\right) \le K(4)\exp(-|x|^2).
\]
This further implies that
\begin{align*}
1 &\le 3 K(4) \exp\left( \frac{|x|^2}{2t+1} - |x|^2 \right)
\\&= 3 K(4) \exp\left( \frac{|x|^2 -2t|x|^2 -|x|^2}{2t+1} \right)
\\& = 3 K(4) \exp\left( \frac{-2t|x|^2 }{2t+1} \right)
\end{align*}
and this must hold for all $t \in [0,4]$ and $x \in \R$. However, as $|x| \to \infty$ then
\[
3 K(4) \exp\left( \frac{-2t|x|^2 }{2t+1} \right) \to 0
\]
and so
\begin{equation} \label{E:contradiction}
1 \le 3 K(4) \exp\left( \frac{-2t|x|^2 }{2t+1} \right)
\end{equation}
cannot hold for all $x \in \R$ which is a contradiction. Thus $\exp(-|x|^2)$ is not admissible.

\subsection{$\rho(x) = \exp(-|x|^b)$}

Suppose that $b\in (1,2)$ and that $\rho(x) = \exp(-|x|^b)$.
\begin{align*}
(G(t,\cdot) * \rho)(x) &= \int_{\R} G(t,x-y)\exp(-|y|^b) \ud y
\\&= \sum_{n=0}^\infty \frac{(-1)^n}{n!} \int_{\R} G(t,x-y) |y|^{nb} \ud y
\\& = \sum_{n=0}^\infty \frac{(-1)^n}{n!} \mathbb{E}(|N|^{nb})
\end{align*}
where above $N$ is a normal random variable with mean $x$ and standard deviation $t$. Note that this can't technically be done because the power series for $\exp(x)$ does not converge uniformly on the whole real line, however doing this seems to work as can be seen in Section \ref{S:SpecCase} below. According to \cite{Winklebauer}, this moment is given by
\[
\mathbb{E}(|N|^{nb}) = (2t)^{nb/2} \frac{\Gamma\left(\frac{nb+1}{2}\right)}{\sqrt{\pi}} \phi\left( -\frac{nb}{2}, \frac{1}{2} ; - \frac{x^2}{2t} \right)
\]
where $\phi$ is the Krummer Confluent Hypergeometric function
\[
_1F_1(\alpha,\gamma;z) = \phi(\alpha,\gamma;z) = \sum_{k=0}^\infty \frac{\alpha^{\bar{k}}}{\gamma^{\bar{k}}} \frac{ z^k }{ k! } \quad \text{and} \quad z^{\bar{k}} = \frac{ \Gamma(z+k) }{ \Gamma(z) }.
\]
Plugging this in above gives us that
\begin{align}\label{E:exp(x^b)}
(G(t,\cdot) * \rho)(x) &= \sum_{n=0}^\infty \frac{(-1)^n}{n!}(2t)^{nb/2} \frac{\Gamma\left(\frac{nb+1}{2}\right)}{\sqrt{\pi}} \phi\left( -\frac{nb}{2}, \frac{1}{2} ; - \frac{x^2}{2t} \right)
\\&= \sum_{n=0}^\infty \frac{(-1)^n}{n!} (2t)^{nb/2} \frac{\Gamma\left(\frac{nb+1}{2}\right)}{\sqrt{\pi}} \sum_{k=0}^\infty \frac{(-nb/2)^{\bar{k}}}{(1/2)^{\bar{k}}} \frac{ (\frac{-x^2}{2t})^k }{ k! }
\end{align}
Note that it seems that \eqref{E:exp(x^b)} is only true when $b \in \mathbb{N}$ (see Section \ref{S:MMC} and equation \eqref{E:BRsumrep}).

\bigskip

\textcolor{magenta}{If it works for $b\in \mathbb{N}$, it should work for fractions as well. I can
show you some numerical evidence that why the case $b>1$ won't work during our meeting. We may not
need to spend too much time here unless we can really prove that $\rho$ with $b>1$ is not
admissible.}

\subsection{Special case when $x=0$ and $b=2$}\label{S:SpecCase}
If we set $b=2$ and $x=0$ in \eqref{E:exp(x^b)} then we should be able to recover \eqref{E:exp(x^2)} with $x$ set to $0$, which is
\[
\frac{1}{\sqrt{2t+1}}.
\]
Letting $b=2$ and $x=2$ in \eqref{E:exp(x^b)} gives us
\begin{equation}\label{E:specialcase}
\sum_{n=0}^\infty \frac{(2n)!(-2t)^n}{4^n(n!)^2}
\end{equation}
and plugging this into Wolfram gives us that
\[
\sum_{n=0}^\infty \frac{(2n)!(-2t)^n}{4^n(n!)^2} = \frac{1}{\sqrt{2t+1}} \quad \text{whenever } |t| \le 1/2.
\]
Another thing to note is that the Borel Regularization of the sum \eqref{E:specialcase} is also equal to $1/\sqrt{2t+1}$ for any $t>0$. This can be easily checked by plugging
\begin{verbatim}
Sum[(-2t)^k(2k)!/(4^k(k!)^2), {k, 0, Infinity}, Regularization -> "Borel"]
\end{verbatim}
into Mathematica.

Its possible that we may be able to use this technique to calculate \eqref{E:exp(x^b)} for more general $x$ and $b$. In fact, what I believe that we may be able to show is that if we fix any $t >0$ then $(G(t,\cdot) * |\cdot|^b)(x)$ will decay much faster than $|x|^b$ as $|x| \to \infty$.

\bigskip

\textcolor{magenta}{I can show you during our meeting this week that the case $b=2$ won't work.}


\subsection{More Mathematica Calculations.}\label{S:MMC}
I am unsure that Borel Regularization can be used. For example, I considered $b=3/2$, $t=1/2$ and $x=0$ and if we plug this into \eqref{E:exp(x^b)} then this gives us
\[
(G(1/2,\cdot) * \rho)(0) = \sum_{n=0}^\infty \frac{(-1)^n}{n!} \frac{\Gamma\left(\frac{n(3/2)+1}{2}\right)}{\sqrt{\pi}} \phi\left( -\frac{(3/2)n}{2}, \frac{1}{2} ; 0 \right)
\]
however we get different answers when using Mathematica to calculate this and also this
\[
\int_{\R} G(1/2,y)\exp\left(-|y|^{(3/2)}\right) \ud y.
\]
The codes used were
\begin{verbatim}
Sum[((-1)^n Gamma[((3/2)*n+1)/2]Hypergeometric1F1[(-(3/2)*n/2),1/2,0])/(Sqrt[Pi]n!),

{n,0,Infinity}, Regularization ->"Borel"]
\end{verbatim}
and
\begin{verbatim}
Integrate[(1/Sqrt[Pi])Exp[-Abs[y]^2]Exp[-Abs[y]^(3/2)],{y,-Infinity,Infinity}].
\end{verbatim}
However, it seems that the following is true for every $b \in \mathbb{N}$
\begin{equation}\label{E:BRsumrep}
\sum_{n=0}^\infty \frac{(-1)^n}{n!} \frac{\Gamma\left(\frac{n(b)+1}{2}\right)}{\sqrt{\pi}} \phi\left( -\frac{(b)n}{2}, \frac{1}{2} ; 0 \right) = 	\int_{\R} G(1/2,y)\exp(-|y|^b) \ud y
\end{equation}
where the above sum is the Borel regularization.

Lastly, what I wished to be able to do what calculate
\[
\int_{\R} G(t,x-y)\exp(-|y|^b) \ud y
\]
on Mathematica for arbitrary $t,x$ and $b$ but I was not able to do so. The reason it gave me was that it took too much computational time for my account so maybe you could do it with your account. However to simplify this, we could just fix a $t$ value, say $t=1/2$ and this simplifies the integral to
\[
\frac{1}{\sqrt{\pi}} \int_{\R} \exp(-|x-y|^2)\exp(-|y|^b) \ud y.
\]
However I still cant calculate this on Mathematica. What I believe to be true is that if we assume by contradiction that for $T=1/2$, there exists some constant $K(1/2)$ such that for any $0 \le t \le 1/2$ and $x \in \R$
\[
\frac{1}{\sqrt{\pi}} \int_{\R} \exp(-|x-y|^2)\exp(-|y|^b) \ud y \le K(1/2)\exp(-|x|^b)
\]
then we will come up with a contradiction. If we use Mathematica, we see that for the case $t=1/2$ and $b=2$ then
\[
\lim_{x \to \infty}	\dfrac{	\frac{1}{\sqrt{\pi}} \int_{\R} \exp(-|x-y|^2)\exp(-|y|^2) \ud y}{\exp(-|x|^2)} = \infty
\]
which again shows that $\exp(-|x|^2)$ is not admissible. But Mathematica would not calculate the corresponding limit for any particular $b\in(1,2)$. For example, I tried taking $b=3/2$, which is the limit
\[
\lim_{x \to \infty}	\dfrac{	\frac{1}{\sqrt{\pi}} \int_{\R} \exp(-|x-y|^2)\exp(-|y|^{3/2}) \ud y}{\exp(-|x|^{3/2})}
\]
and the codes I used were:
\begin{verbatim}
f[x_] := Integrate[(1/Sqrt[Pi])Exp[-Abs[x-y]^2]Exp[-Abs[y]^(3/2)],{y,-Infinity,Infinity}]

Limit[f[x]Exp[Abs[x]^(3/2)],x ->Infinity]
\end{verbatim}
but Mathematica would time out each time I tried.

\bigskip
\textcolor{magenta}{Here you might want to try ``NIntegrate'' instead of ``Integrate''.}

\section{Extending the results to the fractional heat equation.}
The following equation driven by space-time white noise is studied in \cite{CDfracheat15}
	\[
	\begin{cases}
		\left( \frac{\partial}{\partial t} -{ _xD_\delta^a} \right)u(t,x) = \rho(u(t,x)) \dot{W}(t,x), & t>0, \; \text{and } x \in \R
		\\ u(0,\cdot) = \mu(\cdot).
	\end{cases}
	\]
The following moment bounds are proven there:
	\begin{equation}\label{E:fracheatspde} \Norm{u(t,x)}_p^2 \le
	\begin{cases}
	J_0^2(t,x) + ( [ \bar{\varsigma} + J_0^2] \star \bar{\mathcal{K}})(t,x), & p=2
	\\  	2J_0^2(t,x) + ( [ \bar{\varsigma} + 2J_0^2] \star \hat{\mathcal{K}}_p)(t,x), & p>2,
	\end{cases}
	\end{equation}
and $\mathcal{K}$ has the following upper bound
	\[
		\mathcal{K}(t,x) \le \frac{C}{t^{1/a}} {_\delta G_a}(t,x)\left(  1+ t^{1/a} \exp( \gamma^{a^*} t) \right).
	\]
\subsection{Discussion.}
	\begin{enumerate}
		\item We run into the issue with this bound not being bounded in $t$
		\item The space time convolution in the moment bound versus just a spacial convolution may cause problems. When dealing with the heat equation and a noise white in time and colored in space, the moment bound only has a spacial convolution.
		\item We may be able to study equation \eqref{E:fracheatspde} but with a noise that is white in time and colored in space and get a result similar to the heat equation case and use this to show the existence of an invariant measure.
	\end{enumerate}

\section{$\Upsilon(\beta) \vert_{\beta =0} < \infty$ versus the spectral density integrability condition in \cite{TZ98} (hypothesis 2.1).}
Hypothesis 2.1 \cite{TZ98} says that the spectral density $\hat{f}$ must satisfy the property that
	\begin{equation}\label{E:hyp2.1}
 \hat{f} \in L^p(\R^d) \quad \text{for some } p \in (d/(d-2), \infty]
 	\end{equation}
Our assumption is that
	\begin{equation}\label{E:speccond}
		\Upsilon(0) = (2\pi)^{-d} \int_{\R^d} \frac{\hat{f}(\xi)}{|\xi|^2} \ud \xi < \infty, \quad \quad d \ge 3.
	\end{equation}

\subsection{Assuming \eqref{E:speccond}.}

I believe that we may find a $\hat{f}$ that satisfies \eqref{E:speccond} but not \eqref{E:hyp2.1}. First, suppose \eqref{E:speccond} is true. Then we know that $\hat{f}(\xi)$ may have some relaxed decay properties at $|x| = \infty$ because of the help from the $|\xi|^{-2}$ term. In particular we can have decay as slow as $|\xi|^{-a}$ for $a \in (d-2,\infty)$ at infinity. At the same time, the $|\xi|^{-2}$ forces $\hat{f}$ to have growth no faster than $|\xi|^{-a}$ for $a \in (0,d-2)$ around the origin. So, if we assume $\hat{f}$ satisfies the condition \eqref{E:speccond} then $\hat{f}$ must be integrable around the origin but may not be integrable at infinity because $\hat{f}$ may have tails that decay as slow as $|\xi|^{-a}$ for $a \in (d-2,d)$.

Considering the above, it seems as if we assume $\hat{f}$ satisfies \eqref{E:speccond}, then $\hat{f}$ may not be able to satisfy the condition \eqref{E:hyp2.1}. First, let $d \ge 3$ and consider $0 < \delta < K$ where $K$ will be determined below. Now consider $\hat{f}$ such that $\hat{f}$ behaves like $|\xi|^{-(d-2-\delta)}$ around the origin and $|\xi|^{-(d-2+\delta)}$ at infinity. BWOC suppose that $\hat{f}$ satisfies condition \eqref{E:hyp2.1}. Then to take care of the thick tails, we would need $p$ to be such that
	\begin{equation}
		p > \frac{d}{d-2+\delta}
	\end{equation}
but recall that $p > d/(d-2)$ so we must choose a $\gamma >0$ such that
	\[
		\frac{d+\gamma}{d-2+\delta} > \frac{d}{d-2}
	\]
and thus we need
	\begin{equation}\label{E:pcon}
		p \ge \frac{d+\gamma}{d-2+\delta}.
	\end{equation}
However, at the same time we still need to maintain integrability around the origin. But we will now see that if we choose $\delta>0$ small enough that we will lose integrability around the origin. Note that because of \eqref{E:pcon} we may say that
	\[
		p(d-2-\delta) \ge \frac{(d+\gamma)(d-2-\delta)}{d-2+\delta}
	\]
and if we set
	\[
		\frac{(d+\gamma)(d-2-\delta)}{d-2 + \delta} \ge d
	\]
then this is satisfied when
	\begin{equation}\label{E:hyp2.1fail}
		\frac{d-2-\delta}{d-2 + \delta} \ge \frac{d}{d+\gamma}
	\end{equation}
and since $\delta>0$ is arbitrary and
	\[
		\lim_{\delta \to 0} \frac{d-2-\delta}{d-2 + \delta} = 1 > \frac{d}{d + \gamma}
	\]
then there exists a $K>0$, independent of $p$, such that for any $\delta \in (0,K)$ that \eqref{E:hyp2.1fail} is true. And therefore this implies that
	\[
			p(d-2-\delta) \ge d
	\]
and so $\hat{f}$ cannot be integrable around the origin which means that $\hat{f}$ satisfies \eqref{E:speccond} but not \eqref{E:hyp2.1}.

\subsection{Assuming \eqref{E:hyp2.1}.}

I believe that an additional property of the spectral measure, $\hat{f}$, is that $\hat{f}(\xi) < \infty$ if an only if $t \not= 0$. This is assumed in \cite{BC16}.

So lets assume that $\hat{f}$ blows up at the origin. Also assume that $\hat{f}$ satisfies the condition \eqref{E:hyp2.1}. I believe that even if we assume that $\hat{f}$ must blow up at the origin, then we may still construct an example that satisfies \eqref{E:hyp2.1} but not \eqref{E:speccond} which would conclude that not one condition is more general than the other.

Lets consider an $\hat{f}(\xi)$ such that around the origin, $\hat{f}$ behaves like $|\xi|^{-(d-2-\delta)}$ for some arbitrary $\delta >0$. Suppose the tail of $\hat{f}$ behaves like $|\xi|^{-(d-2)}$. We want to find a $p > d/(d-2)$ such that $\hat{f} \in L^2(\R^d)$.  Because of the behavior around the origin, we would need that
	\[
		p < \frac{d}{d-2-\delta}.
	\]
But
	\[
		\frac{d}{d-2-\delta} > \frac{d}{d-2}
	\]
and so we must restrict our choice of $p$ to be such that
	\begin{equation}\label{E:pinterval}
		\frac{d}{d-2} < p < \frac{d}{d-2-\delta} .
	\end{equation}
Notice that because of the tail behavior, that we only need
	\[
		p(d-2) > d
	\]
which means that any $p$ satisfying \eqref{E:pinterval} will make the tail $L^p(\R^d)$ integrable. Therefore, for such an $\hat{f}$ with origin and tail behavior as described above, we may be able to find a $p$ such that $\hat{f} \in L^p(\R^d)$.

But this $\hat{f}$ will not satisfy the condition \eqref{E:speccond} because the tail behavior of $|\xi|^{-2}\hat{f}(\xi)$ will be $|\xi|^{-d}$, and this is not integrable.

\bigskip
\textcolor{magenta}{Your construction/proof is convincing! This second case suggests that
\cite{TZ98} has to contain some flaw! Can you apply my results in \cite{CK19} to show that the case
with
\begin{equation*}
	\widehat{f} (\xi) = \frac{1}{1+|\xi|^{d-2}} , \qquad d\ge 3
\end{equation*}
is a legal spectral density (namely, nonnegative and nonnegative definite), and solution is
intermittent (namely, no phase transition).
}

\section{$\tilde{\Gamma} = \widehat{|\gamma^{1/2}|} *  \widehat{|\gamma^{1/2}|}$  }
Simplifying the function
	\begin{equation}\label{E:GammaTilde}
		\tilde{\Gamma} := \widehat{|\gamma^{1/2}|} *  \widehat{|\gamma^{1/2}|} = \mathcal{F}(|\gamma^{1/2}|) * \mathcal{F}(|\gamma^{1/2}|)
	\end{equation}
gives us up to a constant
	\[
		\tilde{\Gamma}(\xi) = \mathcal{F}(|\gamma|)(\xi)
	\]
Last time we met you mentioned to look at
	\[
		g(x) = \prod_{i=1}^3 1_{[-1,1]}(x_i), \quad \quad \mathcal{F}g(\xi) = 8\prod_{i=1}^3 \frac{\sin(\xi_i)}{\xi_i}.
	\]
Is $g$ supposed to represent $\gamma$? If so then I think this is a problem. Let $f$ be the spatial correlation function. Since $\gamma$ is the spectral density, then up to a constant, we would have
	\[
		f(x) = \mathcal{F}(g)(x) = 8\prod_{i=1}^3 \frac{\sin(x_i)}{x_i}.
	\]
but this is a problem because the correlation function needs to be nonnegative and nonnegative definite and $\sin(x)/x$ is not nonnegative.

I think we made another mistake as well. We defined $q(x) = (g*g)(x)$ and noted that $\mathcal{F}q = (\mathcal{F}g)^2$. However, we are never taking the Fourier transform of $g*g$. \eqref{E:GammaTilde} is a convolution of Fourier transfroms, not a Fourier transform of convolutions. Using \eqref{E:GammaTilde}, we wont get the
	\[
		\left|\frac{\sin(x)}{x}\right|
	\]
term since by using \eqref{E:GammaTilde},
	\[
		\tilde{\Gamma}(x) = \mathcal{F}\left( |g(\cdot)| \right)(x)  = \mathcal{F}\left( g(\cdot)\right)(x) = \prod_{i=1}^3\frac{\sin(x_i)}{x_i}.
	\]
In our meeting we did
	\[
		\sqrt{\mathcal{F}q} = \sqrt{(\mathcal{F}f)^2} = 8\prod_{i=1}^3 \left| \frac{\sin(x_i)}{x_i} \right|
	\]
but I don't think this should be done.

\bigskip

\textcolor{magenta}{Let $f(x) = (g*g)(x)$.  $h$ is a tent function as I drew in the paper during our
meeting. Please show that $h$ is nonnegative and nonnegative definite. Moreover, how would you find
$\sqrt{\mathcal{F}f (\xi)}$?}

Isn't the condition 	
	\begin{equation}\label{E:GamTilCon}
	\tilde{\Gamma} = \mathcal{F}(|\gamma^{1/2}|) * \mathcal{F}(|\gamma^{1/2}|), \quad \quad \frac{\Gamma(d/2 -1 )}{4 \pi^{d/2}} \int_{\R^d} \tilde{\Gamma}(\xi) |\xi|^{d-2} \ud \xi < \infty
	\end{equation}
the same thing as Dalang's condition? I don't think the absolute value adds any restriction because up to a constant we have:
	\begin{align*}
			\tilde{\Gamma} &= \mathcal{F}(|\gamma^{1/2}|) * \mathcal{F}(|\gamma^{1/2}|) 
			\\ &= \mathcal{F}\left( |\gamma^{1/2}| \cdot |\gamma^{1/2}| \right)
			\\& = \mathcal{F}( |\gamma| )
			\\& = \mathcal{F}( \gamma )
			\\& = f
	\end{align*}
where above $f$ is the correlation function and $\gamma$ is the spectral density and hence both are nonnegative and nonnegative definite. Therefore, \eqref{E:GamTilCon} can be written as the following:
	\[
		\int_{\R^d} \tilde{\Gamma}(\xi) |\xi|^{d-2} \ud \xi = 	\int_{\R^d} f(\xi) |\xi|^{d-2} \ud \xi
	\] 
which is the same thing as Dalang's condition. 

\section{Meaning of $\mathcal{F}(\gamma^{1/2})$.}
We discussed last meaning that equation (3.4) \cite{TZ98}, which says 
	\[
		\tilde{\Gamma} := |\widehat{\gamma^{1/2}}|*|\widehat{\gamma^{1/2}}|
	\]
means that we are calculating $\mathcal{F}(\gamma^{1/2})$ and then taking the absolute value and then doing the convolution. So in other words 
	\begin{equation}\label{E:GamTil}
			\tilde{\Gamma} := |\mathcal{F}(\gamma^{1/2})|* |\mathcal{F}(\gamma^{1/2})|.
	\end{equation}
We saw that this was the case by observing the proof of Theorem 3.3 \cite{TZ98}.
	
We then wanted to consider the case where the correlation function, $f$, is given by 
	\[
		f(x) =(2-|x|)1_{[-2,2]}(x)
	\]
which means that 
	\[
		\hat{f}(\xi) = \gamma(\xi) = \frac{\sin^2(x)}{x^2}.
	\]
Because $\gamma(\xi)$ touches $0$ an infinite number of times then we have an infinite number of options to choose from for $\sqrt{\gamma(\xi)}$. Notice that \eqref{E:GamTil} does not tell us which branch of $\sqrt{\gamma}$ to choose since the absolute value is not taken until after taking the Fourier transform. We considered the following two branches for $\sqrt{\gamma}$:
	\begin{equation}\label{E:SqrtGamma}
		\frac{\sin(\xi)}{\xi} \quad \text{or} \quad \frac{|\sin(\xi)|}{|\xi|}
	\end{equation}
however it seems like we could also choose 
	\begin{equation}\label{E:SqrtGammaK}
		\frac{|\sin(\xi)|}{|\xi|} K(\xi)
	\end{equation}
where 
	\[ K(\xi) = 
		\begin{cases}
			-1 & \xi \in (2\pi, 3\pi)
			\\ 1 & \text{otherwise}
		\end{cases}
	\]
and we could come up with an infinite number of choices for $\sqrt{\gamma}$ of the form \eqref{E:SqrtGammaK}. We wanted to choose the second choice from \eqref{E:SqrtGamma} but what tells us that this is the choice we are supposed to choose? It seems to me now that there are an infinite number of possibilities to define $\tilde{\Gamma}$ and note that if we just consider the two choices from \eqref{E:SqrtGamma} then we can get that 
	\[
		\tilde{\Gamma} = \left| \mathcal{F}\left(\frac{\sin(\cdot)}{\cdot} \right) \right|* \left| \mathcal{F}\left(\frac{\sin(\cdot)}{\cdot} \right) \right| 
			\quad	\text{or}	\quad
		\tilde{\Gamma} = \left| \mathcal{F}\left(\left|\frac{\sin(\cdot)}{\cdot} \right|\right) \right|* \left| \mathcal{F}\left(\left|\frac{\sin(\cdot)}{\cdot} \right| \right) \right| 
	\] 
and I do not believe that these two choices for $\tilde{\Gamma}$ are equal. 

This makes me believe that there is something not being said in this paper that tells us which branch of $\sqrt{\gamma}$ needs to be chosen. I think this has something to do with the centered equation right below equation (2.3) \cite{TZ98} which says that 
	\begin{equation}\label{E:SqrtGammaM}
		M(\varphi)\psi = \varphi(\mathcal{F}(\gamma^{1/2}) * \psi)
	\end{equation}
where
	\[
		X^x(t) = S(t)x + \int_0^t S(t-s)M(B(X^x(s))) dW_s
	\]
and 
	\[
		B(\phi)(\xi) = b(\xi,\phi(\xi)).
	\]
I will look through \cite{DZ92} and \cite{DZp92} to see if this has to do with the construction of the DaPrato integral. 

\section{Proof of Theorem 3.1 \cite{TZ98}} 
The proof starts by saying that the result can be proven once we show that for any $\epsilon \in (0,1)$ and $R>0$ that there exists a compact set $\mathscr{K}(\theta) \subset L^2_{\rho}$ such that 
	\begin{equation}\label{E:3.1ineq}
		\mathbb{P}\left\{ X^{\hat x}(t) \in \mathscr{K}(\theta) \right\} \ge (1-\epsilon)\mathbb{P} \left\{ | X^{\hat x}(t-1)|_{\hat \rho} < R \right\} \quad \text{ for all} t>1.
	\end{equation}

In \textbf{step 3} of this proof they define the set 
	\begin{equation}\label{E:KTheta}
		\mathscr{K}(\theta) := \left\{ S(1)y + \frac{\sin(\alpha \pi)}{\pi} F_{\alpha}h  \; : \; |y|_{\hat{\rho}} \le \theta \text{ and } |h|_{L^q(0,1,L^2_{\rho})} \le \theta\right\}
	\end{equation}
and they also mention that by the factorization theorem from \cite{DZ96} that 
	\begin{equation}\label{E:Factorization}
		X^x(1) = S(1)x + \frac{\sin(\alpha \pi)}{\pi} F_{\alpha}\left( Y^x(\cdot) \right).
	\end{equation}
		
They claim that since
	\begin{equation}\label{E:KTheta1}
			\mathbb{P}\left\{ X^x(1) \in \mathscr{K}(\theta)  \right\} \ge (1-k_3\theta^{-q}(1+R^q))
	\end{equation}
then we have 
	\begin{equation}\label{E:KThetaRecursive}
		\mathbb{P}\left\{ X^x(t) \in \mathscr{K}(\theta)  \right\} \ge (1-k_3\theta^{-q}(1+R^q)) \mathbb{P}\left\{ \left| X^x(t-1) \right|_{\hat{\rho}} < R  \right\} \quad \text{for all } t>1.
	\end{equation}
I am not sure why this true. I was thinking that it follows from something like the following:
	\begin{align*}
			\mathbb{P}\left\{ X^x(t) \in \mathscr{K}(\theta)  \right\}  &= 	\mathbb{P}\left\{ X^x(t) \in \mathscr{K}(\theta) \bigg\vert  \left| X^x(t-1) \right|_{\hat{\rho}} < R  \right\} \mathbb{P}\left\{ \left| X^x(t-1) \right|_{\hat{\rho}} < R  \right\} 
			\\ &\quad \quad \quad + \mathbb{P}\left\{ X^x(t) \in \mathscr{K}(\theta) \bigg \vert \left| X^x(t-1) \right|_{\hat{\rho}} \ge R  \right\} \mathbb{P}\left\{ \left| X^x(t-1) \right|_{\hat{\rho}} \ge R  \right\}.
	\end{align*}
However, we would want something like 
	\begin{equation} \label{E:RUpBd}
		R > \left| X^x(t-1) \right|_{\hat{\rho}}
	\end{equation}
 which implies that 
	\[
		 \mathbb{P}\left\{ \left| X^x(t-1) \right|_{\hat{\rho}} \ge R  \right\} = 0
	\]
and so 
	\begin{equation}\label{E:KThetaConditional}
			\mathbb{P}\left\{ X^x(t) \in \mathscr{K}(\theta)  \right\}  = 	\mathbb{P}\left\{ X^x(t) \in \mathscr{K}(\theta) \bigg\vert  \left| X^x(t-1) \right|_{\hat{\rho}} < R  \right\} \mathbb{P}\left\{ \left| X^x(t-1) \right|_{\hat{\rho}} < R  \right\}.
	\end{equation}
This can't be the case though because this would imply that the following function is bounded, which we do not know to be true:
	\[
		(t,\omega) \mapsto \left| X^x(t-1) \right|_{\hat{\rho}}(\omega) \quad \text{for all } 	(t,\omega) \in (1,\infty) \times \Omega.
	\]
However, we do know that 
	\[
		\sup_{t>\alpha} \mathbb{E}\left[ |X^x(t,\cdot)|_{\hat \rho}^2 \right] < \infty 
	\]
so maybe this helps? However the statement in \eqref{E:3.1ineq} says that it needs to be for arbitrary $R>0$ and \eqref{E:RUpBd} implies that $R$ must be taken sufficiently big and thus $R$ is not arbitrary. 

Now just assuming the above is it true, can we say that the SHE is time homogeneous in the sense that 
	\[
			\mathbb{P}\left\{ X^x(t) \in \mathscr{K}(\theta) \bigg\vert  \left| X^x(t-1) \right|_{\hat{\rho}} < R  \right\} = \mathbb{P}\left\{ X^x(1) \in \mathscr{K}(\theta) \bigg\vert  \left| X^x(0) \right|_{\hat{\rho}} < R  \right\}? 
	\]
Because if this is true then because of \eqref{E:KTheta1} we have that 
	\begin{align*}
			\mathbb{P}\left\{ X^x(t) \in \mathscr{K}(\theta) \bigg\vert  \left| X^x(t-1) \right|_{\hat{\rho}} < R  \right\} &= \mathbb{P}\left\{ X^x(1) \in \mathscr{K}(\theta) \bigg\vert  \left| X^x(0) \right|_{\hat{\rho}} < R  \right\}
			\\ & \ge (1-k_3\theta^{-q}(1+R^q))
	\end{align*}
and combining this with \eqref{E:KThetaConditional} gives us \eqref{E:KThetaRecursive}.

\section{Two possible assumptions for the initial condition.}
We have been deciding between which of the following two assumptions we want to assume:
	\begin{enumerate}
		\item 
			\begin{equation}\label{E:A1} 
				\sup_{t>0} \int_{\R^d} (G(t,\cdot) * |\mu|)^2(x)\rho(x) \ud x < \infty
			\end{equation}
		\item 
			\begin{equation} \label{E:A2} 
				\sup_{x \in \R^d} (G(t,\cdot) * |\mu|)(x) < K(t) \quad \text{for some } K(t) \in \mathbb{N}.
			\end{equation}
	\end{enumerate}

You originally suggested Condition \eqref{E:A1} because it gave a nice relation between the initial condition and the weight and that it may have been able to replace the admissible weight condition. However we saw last meeting that we still need to assume the admissible weight condition because of Proposition 2.1 \cite{TZ98}. 

The other issue with Condition \eqref{E:A1} is that 
	\[
		\sup_{t>0} \int_{\R^d} (G(t,\cdot) * \delta_0)^2(x)\rho(x) \ud x = \infty.
	\]
In order to fix this issue with $\delta_0$, we discussed that we may potentially be able to replace Condition \eqref{E:A1} with the following:
	\begin{equation}\label{E:A1.1}
		\sup_{t>\lambda} \int_{\R^d} (G(t,\cdot) * |\mu|)^2(x)\rho(x) \ud x < \infty \quad \text{for any } \lambda > 0.
	\end{equation}

Now if we assume Condition \eqref{E:A1.1}, then for $\varphi \in L^2_{\rho}$, we will only be able to show that
	\[
		\sup_{t>\lambda} \mathbb{E} \left( |u(t,\cdot)|_{\rho} \right) < \infty.
	\]
Using this, we will then only be able to show the following bounded in probability statement: for all $\epsilon >0$ there exits an $R>0$ such that
	\[
		 \mathbb{P} \left( |u(t,\cdot)|_{\rho} \ge R \right) < \epsilon \quad \text{for all } t \ge \lambda.
	\]
I have not verified this, but I believe that if we assume \eqref{E:A1.1} then our version of \eqref{E:3.1ineq} becomes: for $0 < t_0 \le 1$ and $\epsilon \in (0,1)$ and $R>0$, there exists a compact set $\mathscr{K}(\theta) \subset L^2_{\rho}(\R^d)$ such that 
	\begin{equation}\label{E:Conjecture}
			\mathbb{P}\left\{ X^{\hat x}(t) \in \mathscr{K}(\theta) \right\} \ge (1-\epsilon)\mathbb{P} \left\{ | X^{\hat x}(t-t_0)|_{\hat \rho} < R \right\} \quad \text{for all } t> \lambda + t_0
	\end{equation}
If we assume all of this then we can prove that the laws
	\[
		\left\{\mathscr{L}(u(t,\cdot))\right\}_{t>t_0 + \lambda}
	\]
are tight and hence there exists an invariant measure. So we get the same conclusion as long as the conjecture \eqref{E:Conjecture} is correct. 

\section{I do not think \eqref{E:A1.1} will work as an assumption.}

Recall that we defined for $f(s,y) \in L^q(0,1,L^2_{\rho})$
	\[
		\left(F_{\alpha}^{t_0} f\right)(x) := \int_0^{t_0} ds \int_{\R^d} dy \: (t_0-s)^{\alpha-1}G(t_0-s,x-y)f(s,y) .
	\]
and
	\[
		Y^\varphi(s,y) =\int_0^s \int_{\R^d} (s-r)^{-\alpha}G(s-r,y-z)b(u^\varphi(r,z))W(dr,dz).
	\]
In Lemma 4.1, we proved that (there is a typo in our paper, the integral needs to go from $0$ to $t_0$)
	\begin{equation}\label{E:YNorm}
		\mathbb{E}\int_0^{t_0} Y^{\varphi}(s,\cdot)|_{\rho}^q \le |\varphi|_\rho^q.
	\end{equation}
In order to prove this result we either need \eqref{E:A1} or \eqref{E:A2}. If we are only going to assume \eqref{E:A1.1} then all we will be able to say is that
	\begin{equation}\label{E:YNorm1.1}
		\mathbb{E}\int_\lambda^{t_0} Y^{\varphi}(s,\cdot)|_{\rho}^q \le |\varphi|_\rho^q
	\end{equation}
and this will not be enough in order to prove the conjecture \eqref{E:Conjecture}. 

This is because by Lemma 4.3 in our paper we have that 
	\begin{align*}
		\dfrac{\sin(\alpha \pi)}{\pi}\int_0^t(t-s)^{\alpha-1}\left[G(t-s,\cdot)*Y(s,\cdot)\right](x) \ud s &=	\dfrac{\sin(\alpha \pi)}{\pi}\left(F_{\alpha}^{t_0}Y^{\varphi}\right)(x)
		\\&= \int_0^t  \int_{\R^d} G(t-r,x-z)b(u(r,z)) W(\ud r,\ud z)
	\end{align*}
and thus 
	\[
		u^\varphi(t_0,x) = (G(t_0,\cdot) * \varphi)(x) + \dfrac{\sin(\alpha \pi)}{\pi}\left(F_{\alpha}^{t_0}Y^{\varphi}\right)(x)
	\]
and if you examine the proof of Lemma 3.1, this factorization is very important and in particular, the inequality \eqref{E:YNorm}. I do not see how \eqref{E:YNorm1.1} could be used to replace it. 

One last thing to add is that this problem does not present itself with the original assumption \eqref{E:A2} so we may want to stick with that in order to apply the $\delta_0$ initial condition. However, this means we can only used bounded initial conditions. 

\section{Tightness of $U(t_0,\cdot)$}
Define 
	\begin{equation}
	\label{E:U}
		U(T,\cdot)(A) = \frac{1}{T} \int_{t_0}^{T + t_0} \mathscr{L}(u_{u_0}(t,\cdot))(A) \; \ud t
	\end{equation}
where $u_0$ is the initial condition and $A \in \mathscr{B}(L^2_{\rho}(\R^d))$. We see that $\mathscr{L}(u_{u_0}(t,\cdot))(A)$ defines a probability measure on $\mathscr{B}(L^2_{\rho}(\R^d))$. Because $u_{u_0}(t,x)$ is jointly continuous on $(0,\infty) \times \R^d$ then I believe that $t \mapsto \mathscr{L}(u_{u_0}(t,\cdot))(A)$ is continuous. 

First I will show that \eqref{E:U} defines a probability measure. Lets consider a general measure space $(E, \mathscr{E})$ and let $\mathcal{M}_1(E)$ denote the set of probability measures on $(E, \mathscr{E})$. Define 
	\begin{equation}
		\label{E:IntProb}
		P_T(A) := \frac{1}{T-a} \int_a^T p_t(A) \; \ud t,
	\end{equation}
where $T>a \ge 0$, $t \in (a,T)$, $p_t \in \mathcal{M}_1(E)$, and  $t \mapsto p_t(A)$ is continuous for each $A \in \mathscr{E}$.

\noindent \textbf{\eqref{E:IntProb} defines a Probability Measure:} First, consider the probability measures $p_1,\cdots, p_n \in \mathcal{M}_1(E)$. It is easy to verify that $(p_1 + \cdots + p_n)/n \in \mathcal{M}_1(E)$. Now suppose that for $T > a \ge 0$ and $t \in (a,T)$ that $p_t \in \mathcal{M}_1(E)$ and that $t \mapsto p_t(A)$ is continuous for each $A \in \mathscr{E}$. Then the Lebesgue integral in \eqref{E:IntProb} exists for each $A \in \mathscr{E}$. It remains to show that $P_T \in \mathcal{M}_1(E)$. Note that 
	\begin{align*}
		P_T(A) = \lim_{t_n \in P_n} \frac{1}{T-a}\sum_{t_k \in P_n} p_{t_k}(A) \Delta_n
	\end{align*} 
where $\Delta_n = (T-a)/n$ and $P_n$ is a partition of $[a,b]$ constructed from such $\Delta_n$. From this we easily deduce that 
	\[
		P_T(E) = \lim_{t_n \in P_n} \frac{1}{T-a}\sum_{t_k \in P_n} p_{t_k}(E) \Delta_n = \lim_{t_n \in P_n} \frac{1}{T-a}\sum_{t_k \in P_n} \Delta_n = 1. 
	\]
In addition, it is clear that $P_T(A) \ge 0$ so it remains to show countable additivity. Suppose $\{A_j\}_{j \ge 1} \in \mathscr{E}$ are disjoint. Then,
	\begin{align*}
		P_T\left(\sum_{j \ge 1} A_j \right) &= \lim_{t_n \in P_n} \frac{1}{T-a}\sum_{t_k \in P_n} p_{t_k}\left(\sum_{j \ge 1} A_j \right) \Delta_n \\
		&= \lim_{t_n \in P_n} \frac{1}{T-a} \sum_{j \ge 1} \sum_{t_k \in P_n} p_{t_k}\left( A_j \right) \Delta_n 
	\end{align*}
where the last line follows because each $p_{t_k} \in \mathcal{M}_1(E)$. Since 
	\[
		\sum_{t_k \in P_n} p_{t_k}\left( A_j \right) \Delta_n \le \sum_{t_k \in P_n}  \Delta_n = 1,
	\]
then the dominated convergence theorem implies that
	\begin{align*}
		\lim_{t_n \in P_n} \frac{1}{T-a} \sum_{j \ge 1} \sum_{t_k \in P_n} p_{t_k}\left( A_j \right) \Delta_n &=  \sum_{j \ge 1} \lim_{t_n \in P_n} \frac{1}{T-a} \sum_{t_k \in P_n} p_{t_k}\left( A_j \right) \Delta_n \\
		&=  \sum_{j \ge 1} p(A_j),
	\end{align*}
and so $P_T$ is countably additive which means $P_T \in \mathcal{M}_1(E)$. 

\begin{remark}
	\normalfont
	It seems that the assumption $t \mapsto p_t(A)$ defines a continuous mapping may be unnecessary. It seems all we would really need is the existence of the Lebesgue integral in \eqref{E:IntProb} which only requires $t \mapsto p_t(A)$ be measurable and bounded almost everywhere. DaPrato \cite{DZ96} requires that $P_t(x,A) = u_x(t,\cdot)(A)$ be a stochastically continuous Feller semi-group in order for \eqref{E:U} to be a probability measure, which happens to be the case in our situation. 
\end{remark}

\noindent \textbf{\eqref{E:IntProb} is tight if the $p_t$ are tight:} Now lets further assume that the family of probability measures $\{p_t\}_{t \in [a,T]}$ is tight. Let $\epsilon>0$ be given. Then there exists a a compact set $K \in \mathscr{E}$ such that for each $t \in [a,T]$, $p_t(K^c) < \epsilon$ where $K^c = E \setminus K$. Then, 
	\begin{align*}
		P_T(K) &= \frac{1}{T-a} \int_a^b p_t(K) \; \ud t \\
		&= \lim_{t_n \in P_n} \frac{1}{T-a}\sum_{t_k \in P_n} p_{t_k}(K) \Delta_n \\
		& \le \epsilon  \lim_{t_n \in P_n} \frac{1}{T-a}\sum_{t_k \in P_n} \Delta_n \\
		&= \epsilon,
	\end{align*}
which implies that the family  $\{P_T\}_{T \ge a}$ defined in \eqref{E:IntProb} is tight.
\begin{small}% {{{ Reference
	\addcontentsline{toc}{section}{Bibliography}
	\begin{thebibliography}{alpha}
		\bibitem{BC16}% {{{
		Balan, R. and Chen, Le
		\newblock Parabolic Anderson Model with Space-Time Homogeneous Gaussian Noise and Rough Initial Condition
		\newblock \textit{J Theor Probab.} (2016)
		% }}}
		\bibitem{CDfracheat15}% {{{
		Chen, Le and Dalang, R.
		\newblock Moments, intermittency and growth indices for the
		nonlinear fractional stochastic heat equation
		\newblock \textit{Stoch PDE: Anal Comp} (2015) 3:360–397
		% }}}
		\bibitem[CK19]{CK19}% {{{
		Chen, Le and Kunwoo Kim.
		\newblock Nonlinear stochastic heat equation driven by spatially colored noise: moments and intermittency.
		\newblock {\em Acta Math. Sci. Ser. B}, 39B (2019), No. 3, 645--668.
		% }}}
		\bibitem[DZ92a]{DZ92}% {{{
		Da Prato, Giuseppe and Jerzy Zabczyk.
		\newblock {\it Stochastic equations in infinite dimensions.}
		Encyclopedia of Mathematics and its Applications, 44. Cambridge University Press, Cambridge, 1992. xviii+454 pp.
		% }}}
		\bibitem[DZ96]{DZ96}% {{{
		Da Prato, Giuseppe and Jerzy Zabczyk.
		\newblock {\it Ergodicity for infinite-dimensional systems. }
		\newblock London Mathematical Society Lecture Note Series, 229. Cambridge University Press, Cambridge, 1996.
		% }}}
		\bibitem[DZ92b]{DZp92}% {{{
		Da Prato, Giuseppe and Jerzy Zabczyk.
		\newblock {\it Invariant measures for semilinear stochastic equations.}
		Stochastic Anal. Appl. 10, 387-408.
		% }}}
		\bibitem[DZ96]{DZ96}% {{{
		Da Prato, Giuseppe and Jerzy Zabczyk.
		\newblock {\it Ergodicity for infinite-dimensional systems. }
		\newblock London Mathematical Society Lecture Note Series, 229. Cambridge University Press, Cambridge, 1996.
		% }}}
		\bibitem[TZ98]{TZ98}% {{{
		Tessitore, Gianmario and Jerzy Zabczyk.
		\newblock Invariant measures for stochastic heat equations.
		\newblock {\it Probab. Math. Statist.} 18 (1998), no. 2, Acta Univ. Wratislav. No. 2111, 271--287.
		% }}}
		\bibitem{Winklebauer}% {{{
		Winklebauer.
		\newblock Moments and Absolute Moments of the
		Normal Distribution.
		\newblock {\it Arxiv } https://arxiv.org/abs/1209.4340.
		% }}}
	\end{thebibliography}
\end{small}
% }}}



\end{document}



